\documentclass[11pt,a4paper]{article}
\usepackage[T1]{fontenc}
\usepackage[utf8]{inputenc}
\usepackage{amsmath,amssymb,xcolor}
\usepackage[margin=25mm]{geometry}
\usepackage[colorlinks=true,urlcolor=blue]{hyperref}
\definecolor{nexorablue}{HTML}{2459B5}
\setlength{\parindent}{0pt}
\setlength{\parskip}{7pt}
\setlength{\emergencystretch}{2em}
\renewcommand{\familydefault}{\sfdefault}
\begin{document}

{\small\bfseries\color{nexorablue} NEXORA CPGE}\par
{\small MPSI / NEX-MES-05 / Corrigé}\par
{\LARGE\bfseries Régression d’une position en fonction du temps}\par
\vspace{0.25em}\hrule\vspace{0.6em}

\section*{Calcul des coefficients}
La somme des produits centrés vaut $9{,}90$. Ainsi $a=1{,}98\,\mathrm{cm\,s^{-1}}$ et $b=5{,}05-1{,}98\times2{,}5=0{,}10\,\mathrm{cm}$. La pente représente la vitesse pour ce modèle.

\section*{Résidus}
Les prédictions sont $2{,}08$, $4{,}06$, $6{,}04$ et $8{,}02\,\mathrm{cm}$. Les résidus valent $0{,}02$, $-0{,}06$, $0{,}06$ et $-0{,}02\,\mathrm{cm}$. Leur somme est nulle à l’arrondi près. Pour juger l’accord quantitatif, il faudrait les comparer aux incertitudes des observations.


\end{document}
